\Askhsh{18562}{4}
\begin{ylh}{1}
    \item 9.2. Το Πυθαγόρειο θεώρημα
    \item 10.2. Εμβαδόν ευθύγραμμου σχήματος - Ισοδύναμα ευθύγραμμα σχήματα.
\end{ylh}
Δίνεται τετράγωνο ΑΒΓΔ πλευράς $\alpha$.
\begin{alist}
    \item Να αποδείξετε ότι το μήκος της διαγωνίου ΒΔ ισούται με $\alpha\sqrt{2}$ και να βρείτε το εμβαδό του. \monades{05}
    \item
    \begin{rlist}
        \item Να σχεδιάσετε το τετράγωνο ΒΔΖΗ έτσι ώστε το σημείο Α να είναι εσωτερικό σημείο του και να αποδείξετε ότι το σημείο Α είναι το κέντρο του τετραγώνου ΒΔΖΗ. \monades{07}
        \item Να υπολογίσετε το εμβαδό του ΒΔΖΗ και να το συγκρίνετε με το εμβαδό του αρχικού τετραγώνου ΑΒΓΔ. \monades{08}
    \end{rlist}
    \item Επαναλαμβάνουμε το σχεδιασμό όπως περιγράφηκε παραπάνω και σχηματίζουμε νέο τετράγωνο, με πλευρά κάθε φορά, τη διαγώνιο του προηγούμενου τετραγώνου. Δηλαδή με πλευρά τη διαγώνιο ΔΗ του τετραγώνου ΒΔΖΗ σχεδιάζουμε νέο τετράγωνο, το ΔΗΘΚ. Με πλευρά τη διαγώνιο ΗΚ του ΔΗΘΚ σχεδιάζουμε νέο τετράγωνο κ.ο.κ. Αν θέλουμε να σχεδιάσουμε τετράγωνο του οποίου το εμβαδό του θα είναι 16 φορές το εμβαδό του αρχικού τετραγώνου ΑΒΓΔ, πόσες φορές ακόμη πρέπει να επαναλάβουμε το σχεδιασμό αυτό; Να δικαιολογήσετε την απάντησή σας. \monades{05}
\end{alist}